% Reference entries, F (AN-11).  Law 1: os/iword/app/face_a.c take_name
% (DOS 8.3 names: letters, digits, _ - and one period; upper case), name_key
% (the offered name, a typed character replaces it, Backspace), M_FONT
% (IW_FONT_MENU; 1 Size, 2 Appearance, 3 Normal; 4 and 5 return to the
% document), codes.c iw_ed_normal, format_key / paint_format (IW_FMT_INIT,
% IW_FMT_LINE_INIT: the screens, their values at the cursor, the
% selections that act).  Shots: ../wordperfect51-decomp mint MINT06, MINT21,
% MINT28.
% GATES: test-pubs-iword test-iword-face test-iword-codes test-flair-iword-codes test-flair-iword-format test-pubs-flair-screens
\Gates{test-pubs-iword test-iword-face test-iword-codes test-flair-iword-codes test-flair-iword-format test-pubs-flair-screens}
\begin{Feature}{Filenames}
Each document you save is a file on the data volume, and you give it a
name when you first save it.  InitechWord uses the file names of
InitechDOS.

A name has one to eight characters, then, if you like, a period and one to
three more characters, for example \Scr{MEMO}, \Scr{REPORT92} or
\Scr{MINUTES.MAR}.  The characters may be letters, digits, the hyphen
(\Scr{-}) and the underscore (\Scr{\_}).  You may type the letters in
capitals or small letters; InitechWord shows and saves them in capitals.
InitechWord does not add a period or an extension of its own: a document
saved as \Scr{MEMO} is the file \Scr{MEMO}.

To type a name at a prompt,
\begin{WPSteps}
\item Type the name.  If the prompt already shows a name, the first
character you type replaces it.
\item Press \K{Backspace} to erase the last character, if you need to.
\item Press \K{Enter}.
\end{WPSteps}

\Notes
\Topic{Names That Are Not Allowed}
An empty name, a name with any other character, with more than eight
characters before the period or three after it, with nothing before the
period, or with two periods is refused with \Scr{ERROR: Invalid file name}.  Press any key to
return to the prompt and correct the name.

\Topic{Where Documents Are Kept}
Documents are saved in the top level of the data volume, which InitechWord
shows as drive \Scr{A:}: a document named \Scr{MEMO} is shown on the
status line as \Scr{A:\wbs MEMO}.  You type only the name, never a drive or
a folder.

\SeeAlso{Retrieve; Save}
\end{Feature}

\begin{Feature}{Font}
The Font key sets how text looks: its size, its appearance, or a return to
normal text.  Its menu appears on the status line.

To use the Font menu,
\begin{WPSteps}
\item Press \FK{Font}{Ctrl-F8}.  The status line shows
\StatusLine{1 Size; 2 Appearance; 3 Normal; 4 Base Font; 5 Print Color: 0}
\item Select \Ch{Size}{1} to choose a size (see \InRef{Size})
\Or
Select \Ch{Appearance}{2} to choose an appearance (see \InRef{Appearance})
\Or
Select \Ch{Normal}{3} to end every size and appearance at the cursor.
\end{WPSteps}
\Documents{wpkey}{Ctrl-F8}
\Documents{wpmenu}{Ctrl-F8,3=Normal}
\FlairShot{font}{640}{440}{The Font menu on the status line}

\Notes
\Topic{Base Font and Print Color}
Selecting Base Font (4) or Print Color (5) returns you to the document
without a change, as does any key that is not on the menu.

\Topic{Block}
With Block on, the Font key shows \Scr{Attribute: 1 Size; 2 Appearance: 0}
and gives the size or appearance to the block (see \InRef{Block,
Attributes}).

\Topic{Normal}
Normal moves the cursor past any off codes right at the cursor; then, for
every size or appearance still in effect at the cursor, it ends it there by
moving its off code to the cursor, and leaves the cursor after the codes.
What you type next is normal text.

\BeginReplay{ref-normal}
For example, press \Press{Bold}{F6} and \Press{Underline}{F8}, type
\TypeRun{both}, then press \Press{Font}{Ctrl-F8}, select \Choose{Normal}{3}
and type \TypeRun{ plain}.  In Reveal Codes (\Keys{Alt-F3}\K{Alt-F3}) the
codes are:
\WordScreen[rows=13]{ref-normal}{Normal: the text after it is plain}
\EndReplay

\SeeAlso{Appearance; Block, Attributes; Size}
\end{Feature}

\begin{Feature}{Format}
The Format key leads to the settings that shape lines of text: the
justification, the line spacing and the left and right margins.  Format
shows a screen of its own instead of the document.

To use Format,
\begin{WPSteps}
\item Press \FK{Format}{Shift-F8}.  The Format screen appears.
\item Select \Ch{Line}{1}.  The Format: Line screen appears.  It shows the
settings in effect at the cursor.
\item Select \Ch{Justification}{3}, \Ch{Line Spacing}{6} or
\Ch{Margins}{7} and make the setting (see \InRef{Justification},
\InRef{Line Spacing} and \InRef{Margins, Left and Right}).
\item Press \FK{Exit}{F7} to return to the document.
\end{WPSteps}
\Documents{wpkey}{Shift-F8}
\Documents{wpmenu}{Shift-F8,1=Line}
\FlairShot{format}{640}{440}{The Format screen}
\FlairShot{formatline}{640}{440}{The Format: Line screen, with the settings InitechWord starts with}

\Notes
\Topic{Leaving the Format Screens}
From Format: Line, \K{F7} returns to the document; \K{Enter}, \Typed{0},
\K{F1} or \K{Esc} return to the Format screen.  From the Format screen,
\K{F7}, \K{Enter}, \Typed{0}, \K{F1} or \K{Esc} return to the document.

\Topic{Other Selections}
The Format screen lists four groups -- Line, Page, Document and Other --
and Format: Line nine settings, as the screens of WordPerfect 5.1 do.
InitechWord acts on Line (1) and, on the Format: Line screen, on
Justification (3), Line Spacing (6) and Margins (7).  Pressing the number of
any other selection leaves the screen as it is.  The values shown beside
those selections -- Hyphenation \Scr{No}, Line Height \Scr{Auto}, Tab Set
\Scr{Rel; -1", every 0.5"} and the others -- are those of a new document,
and InitechWord does not change them.

\Topic{Settings in Effect}
Each setting you make goes into the document as a code at the cursor (see
\InRef{Codes}).  Text before the code is not affected.  The Format: Line
screen always shows the settings in effect where the cursor is.

\Topic{Several Settings at Once}
You can make several settings in one visit to the Format: Line screen:
each goes into the document in the order you make it.

\SeeAlso{Codes; Justification; Line Spacing; Margins, Left and Right}
\end{Feature}
